\documentclass{beamer}

\usepackage{beamerthemesplit}

\title[Finite generation of pluricanonical rings]{Finite generation of pluricanonical rings}
\author{Christopher Hacon}
\institute{University of Utah}
\date{January, 2007}

\usepackage{amsmath,amssymb,amsthm,amscd,epic,eepic}
%\usepackage{amsmath,amssymb,amstext,amsthm,amscd,epic,eepic}

% Theorems
\theoremstyle{plain}
\newtheorem*{thm}{Theorem}
\newtheorem*{pro}{Proposition}
\newtheorem*{lem}{Lemma}
\newtheorem*{cor}{Corollary}
\newtheorem*{con}{Conjecture}
\newtheorem*{que}{Question}
\newtheorem*{cla}{Claim}
\theoremstyle{definition}
\newtheorem*{defi}{Definition}
\newtheorem*{ack}{Acknowledgements}
\theoremstyle{remark}
\newtheorem*{rem}{Remark}
\newtheorem*{rems}{Remarks}
\newtheorem*{ex}{Example}
\newtheorem*{exs}{Examples}

% Varie
\def\rup#1.{\ulcorner#1\urcorner}
\def\R{\mathbb {R}}
\def\Q{\mathbb {Q}}
\def\OO{\mathcal {O}}
\def\bs{\boldsymbol}
\def\ds{\displaystyle}
\def\pf{{\em Proof.}\ \,}
\def\ipf{{\em Idea of the Proof.}\ \,}
\def\br{\buildrel}
\def\ov{\overline}
\def\su{\subset}
\def\suq{\subseteq}
\def\({\left(}
\def\){\right)}
\def\un{\underline}

% Frecce
\def\inc{\hookrightarrow}
\def\fun{\rightarrow}
\def\lfun{{\longrightarrow}}
\def\funn{\twoheadrightarrow}

% Greek
\def\om{\omega}
\def\Th{\Theta}

% \mathbb
\def\bA{{\mathbb{A}}}
\def\bC{{\mathbb{C}}}
\def\bN{{\mathbb{N}}}
\def\bP{{\mathbb{P}}}
\def\bQ{{\mathbb{Q}}}
\def\bR{{\mathbb{R}}}
\def\bZ{{\mathbb{Z}}}
\def\Mob{{\operatorname{Mov}}}
% \mathcal
\def\cA{{\mathcal A}}
\def\cB{{\mathcal B}}
\def\BL{{\mathcal BL}}
\def\cC{{\mathcal C}}
\def\cD{{\mathcal D}}
\def\cE{{\mathcal E}}
\def\cExt{{\mathcal E}xt}
\def\cF{{\mathcal F}}
\def\cG{{\mathcal G}}
\def\cH{{\mathcal H}}
\def\cHom{{\mathcal H}om}
\def\cI{{\mathcal I}}
\def\cL{{\mathcal L}}
\def\cM{{\mathcal M}}
\def\cN{{\mathcal N}}
\def\cO{{\mathcal O}}
\def\cP{{\mathcal P}}
\def\cT{{\mathcal T}}
\def\cTor{{\mathcal T}or}
\def\cU{{\mathcal U}}
\def\cZ{{\mathcal Z}}

% \widetilde
\def\wt{\widetilde}
\def\tC{{\widetilde{C}}}
\def\tE{{\widetilde{E}}}
\def\tH{{\widetilde{H}}}
\def\tX{{\widetilde{X}}}
\def\tY{{\widetilde{Y}}}
\def\tp{{\widetilde{p}}}
\def\tg{{\widetilde{g}}}
\def\tTh{{\widetilde{\Theta}}}
\def\tcA{{\widetilde{\cA}}}
\def\tcM{{\widetilde{\cM}}}
\def\tcU{{\widetilde{\cU}}}
\def\tPic{{\widetilde{\textrm{Pic}}}}

% \mathfrak
\def\mP{\mathfrak{P}}

% Math Text
\def\ch{{\textrm{ch}}}
\def\codim{{\textrm{codim\,}}}
\def\coker{{\textrm{coker\,}}}
\def\ecoker{{\emph{coker}\,}}
\def\Ext{{\textrm{Ext\,}}}
\def\eExt{{\emph{Ext}\,}}
\def\Gr{{\textrm{Gr}}}
\def\Hilb{{\textrm{Hilb\,}}}
\def\eHilb{{\emph{Hilb\,}}}
\def\Hom{{\textrm{Hom\,}}}
\def\eHom{{\emph{Hom}\,}}
\def\id{{\textrm{id}}}
\def\Im{{\textrm{Im\,}}}
\def\im{{\textrm{im\,}}}
\def\NS{{\textrm{NS}}}
\def\Pic{{\textrm{Pic}}}
\def\ePic{{\emph{Pic}}}
\def\etPic{{\widetilde{\emph{Pic}}}}
\def\Re{{\textrm{Re\,}}}
\def\rk{{\textrm{\,rk}}}
\def\Span{{\textrm{Span\,}}}
\def\Spec{{\textrm{Spec\,}}}
\def\Stab{{\textrm{Stab}}}
\def\eStab{{\emph{Stab}}}
\def\td{{\textrm{td}}}
\def\Tor{{\textrm{Tor\,}}}
\def\eTor{{\emph{Tor\,}}}
\def\Tors{{\textrm{Tors\,}}}
\def\eTors{{\emph{Tors}\,}}

\mode<presentation>
{
  \setbeamertemplate{background canvas}[vertical shading][bottom=red!10,top=blue!10]

  \usetheme{Warsaw}
  \usefonttheme[onlysmall]{structurebold}
}

\begin{document}

\begin{frame}
  \titlepage
\end{frame}



\begin{frame}
  \frametitle{Outline of the talk}
  \tableofcontents[pausesections, hideallsubsections]
\end{frame}

\AtBeginSection[]
{
  \begin{frame}
    \frametitle{Outline of the talk}
    \tableofcontents[currentsection, hideothersubsections]
  \end{frame}
}

\section{Introduction}

%\subsection{The minimal model program for 3-folds}

\subsection{The main Theorem}
\begin{frame}
  \frametitle{Notation}
\begin{itemize}
\item I am interested in the birational classification of complex projective
varieties.
\pause
\item $X\subset \mathbb{P}_{\mathbb{C}}^N$ is defined by homogeneous
polynomials $P_1,\cdots , P_t\in \mathbb{C}[x_0,\cdots ,x_N]$.
\pause
\item $X,X'$ are birational if they have isomorphic open subsets $U\cong U'$.
\pause
\item If $X$ is smooth of dimension $n$, then $T_X$ denotes the {\bf tangent bundle} of $X$ and $\omega _X=\Lambda ^nT_X^\vee$ is the {\bf canonical line bundle} of $X$.
\pause
\item $K_X$ denotes (a choice of) the canonical divisor i.e. a formal linear combination of codimension $1$ subvarieties of $X$ given by the zeroes of
a section $s\ne 0$ of $\omega _X$.
\pause
\item We would like to use $\omega _X$ to understand the geometry of $X$.
\end{itemize}

\end{frame}

\begin{frame}
  \frametitle{Birational invariants}
\begin{itemize}
\item If $X$ and $X'$ are birational, then $\forall m\geq 0$
$$H^0(X,\omega _{X}^{\otimes m})\cong H^0(X',\omega _{X'}^{\otimes m}).$$ 
\pause
\item The numbers $P_m(X):=h^0(\omega _{X}^{\otimes m})$ are the {\bf plurigenera} of $X$.
\pause
\item The {\bf pluricanonical ring} of $X$ is given by $$\mathfrak{R}(K_X)=\bigoplus _{m\in \mathbb {N}}H^0(X,\omega _{X}^{\otimes m}).$$ 
\pause
\item The {\bf Kodaira dimension} of $X$ is $$\kappa (X):={\rm tr.deg.}_{\mathbb {C}} \mathfrak{R}(K_X) -1.$$
\pause
\item If $\kappa (X)={\rm dim}X$, we say that $X$ is of {\bf general type}.
\pause
\item One would like to understand the structure of this ring and to use its features to classify complex projective varieties.
\end{itemize}

\end{frame}
\begin{frame}
  \frametitle{Finite generation}
Today I would like to discuss a recent result of
Birkar-Cascini-Hacon-M\textsuperscript{c}Kernan and Siu:
\begin{theorem} Let $X$ be a smooth projective variety, then $  \mathfrak{R}(K_X)$ is finitely generated.\end{theorem}
\pause
The two proofs are independent. Our proof is algebraic and uses the ideas of the MMP (minimal model program). Siu's proof is analytic and requires $X$ to be of general type.
\end{frame}
\subsection{Curves}

\begin{frame}
  \frametitle{Curves}
\begin{itemize}
\item Any two birational curves are isomorphic.
\pause
\item Curves are topologically determined by their genus. 
\pause
\item $g=0$: ${\mathbb P} ^1$ and $\mathfrak{R}(K)\cong \mathbb {C}$
\pause
\item $g=1$: There is a $1$-dimensional family of elliptic curves; $K=0$ so 
$\mathfrak{R}(K)\cong \mathbb {C}[t]$.
\pause
\item $g\geq 2$: Curves of general type. 

These form a $3g-3$-dimensional family.  

$K$ is ample and $|3K|$ determines an embedding.

\pause
\item $\mathfrak{R}(K)$ is then finitely generated (by Serre-Vanishing).
\end{itemize}

\end{frame}

\subsection{Surfaces}

\begin{frame}
  \frametitle{Surfaces of general type}
\begin{itemize}
\item If $S$ is a surface of general type and $K_S$ is nef (i.e. $K_S\cdot C\geq 0$ for all curves $C\subset S$), then we say that $S$ is a minimal surface.
In this case $K_S$ is semiample.
\pause
\item Therefore, there is a birational morphism $\nu :S\to \bar{S}\subset \mathbb{P}^N$ to a surface (with rational double points) so that $\omega _X^{\otimes m}=\nu ^* \OO _{\bar{S}}(1)$.
\pause
\item It is known that $\mathfrak{R}(K_S)$ is finitely generated if and only if $$\mathfrak{R}(K_S)^{(m)}=\bigoplus _{i\in \mathbb N}H^0(\OO _S(imK_S))$$
\vskip-.2cm is finitely generated (for any $m>0$).
\pause
\item By Serre vanishing $\mathfrak{R}(K_S)^{(m)}\cong \mathfrak{R} (\OO _{\bar{S}}(1))$ is finitely generated.
\pause
\item Note that we may take $m=5$, $\nu$ contracts rational curves $E$ with $E\cdot K_S=0$ and $\bar{S}={\rm Proj} \mathfrak{R}(K_S)$. 
\end{itemize}

\end{frame}



\begin{frame}
  \frametitle{Surfaces of general type II}

\begin{itemize}\item If $S$ is of general type, but $K_S$ is not nef, then there is a rational curve $E$ such that $E^2=E\cdot K_S=-1$. 
\pause
\item We may contract $E$, i.e. there is a morphism of smooth surfaces
$\nu_1:S\to S_1$ which is an isomorphism on $S-E$ and such that $E$ maps to a point. 
\pause
\item Since $K_S=\nu_1^*K_{S_1}+E$, we have $\mathfrak{R}(K_S)\cong
\mathfrak{R}(K_{S_1})$.
\pause
\item If $K_{S_1}$ is not nef, we may repeat the procedure and we obtain a sequence of morphisms $S\to S_1\to S_2 \to \cdots$.
\pause
\item It is easy to see that $0<\rho (S_{i+1})=\rho (S_{i})-1$
and so this sequence of morphism is finite, i.e. there is an integer $t>0$ such that $K_{S_t}$ is nef i.e. $S_t$ is a minimal model of $S$.
\pause
\item Then $\mathfrak{R}(K_{S_t})$ is finitely generated and so is 
$\mathfrak{R}(K_S)$.
\end{itemize}

\end{frame}



\begin{frame}
 \frametitle{Surfaces with $\kappa \leq 2$}

\begin{itemize}
\item If $S$ is not of general type, then $\kappa (S)\in \{-1,0,1\}$.
\pause
\item If $\kappa (S)=-1$, then $\mathfrak{R}(K_S)\cong \mathbb{C}$.
\pause
\item If $\kappa (S)=0$, then $\mathfrak{R}(K_S)^{(m)}\cong \mathbb{C}[t]$
for some $m>0$. (For surfaces $m=12$.)
\pause
\item If $\kappa (S)=1$, then we first contract all $-1$ curves to obtain
a morphism $\nu :S\to S'$ such that $K_{S'}$ is nef (i.e. $S'$ is minimal).
Then $K_{S'}$ is semiample and it gives a morphism $f:S'\to C$
whose general fiber is an elliptic curve.
\pause
\item We have that $mK_{S'}=f^*H$ where $m>0$ is an integer and
$H$ is an ample divisor on $C$.
\pause
\item Therefore, $ \mathfrak{R}(K_S)^{(m)}\cong \mathfrak{R}(K_{S'})^{(m)}
\cong \mathfrak{R}( H)$ is finitely generated.
\end{itemize}
\end{frame}

\section{Higher dimensional varieties}
\subsection{Reduction to log pairs of general type}
\begin{frame}
 \frametitle{Log pairs}

\begin{itemize}
\item One would like to generalize the above approach to 
higher dimensional varieties.
\pause
\item In order to achieve this one has to allow varieties with mild singularities and the minimal models will not be unique.
\pause
\item We would also like to have more flexibility and so we work
with log-pairs.
\pause
\item The proof will in fact proceed by induction on the dimension.
We will need to relate $K_X$ to it's restriction to a divisor $S$ (which typically belongs to $S\in |mK_X|$).
\pause
\item This is achieved by the adjunction formula which has the form
$(K_X+S)|_S=K_S+\Delta _S$ for some $\Q$-divisor $\Delta _S$. 
\end{itemize}
\end{frame}
\begin{frame}
 \frametitle{KLT Log Pairs}
\begin{itemize}
\item We consider log pairs $(X,\Delta )$ where $X$ is a normal $\Q$-factorial variety and $\Delta$ is a divisor with positive real coefficients.
\pause
\item Then $K_X+\Delta$ is $\mathbb{R}$-Cartier and so $(K_X+\Delta )\cdot C$ and
$f^*(K_X+\Delta )$ still make sense.
\pause
\item Consider a map $f:X'\to X$ such that $X'$ is smooth and $f^* \Delta +{\rm Exc}(f)$ has simple normal crossings support.
Write $$K_{X'}+\Delta '=f^*(K_X+\Delta ).$$
\pause
\item $(X,\Delta )$ is KLT if each coefficient of $\Delta '$ is less than $1$.
\end{itemize}
\end{frame}
\begin{frame}
 \frametitle{A Theorem of Fujino-Mori}
Mori and Fujino have shown that:
\begin{theorem} Let $(X,\Delta )$ be any KLT $\Q$-factorial pair with 
$\Delta \in {\rm Div }_\Q (X)$. Then there exists a {\bf big} KLT $\Q$-factorial 
pair $(Y,\Gamma )$ with $\Gamma \in {\rm Div }_\Q (Y)$ such that
$$\mathfrak{R}(K_X+\Delta )^{(m)}\cong \mathfrak{R}(K_Y+\Gamma )^{(m)}$$
for any sufficiently divisible integer $m>0$.
\end{theorem}
\pause
Recall that $K_Y+\Gamma$ is big if ${\rm tr.deg.}_{\mathbb C}\mathfrak{R}(K_Y+\Gamma )^{(m)}=\dim Y$ or equivalently if $K_Y+\Gamma\sim _\Q A+B$ with $A$ ample and $B\geq 0$.

\pause
It follows that to show the main theorem (i.e. that $\mathfrak{R}(K_X)$ 
is finitely generated), it suffices to prove the corresponding statement 
for big KLT log pairs.
\end{frame}
\begin{frame}
\subsection{The MMP with scaling}
\frametitle{The MMP with scaling}
\begin{itemize}\item From now on we will consider 
$\Q$-factorial KLT pairs $(X,\Delta )$ 
such that $\Delta$ is big. We allow $\Delta \in {\rm Div }_{\mathbb {R}}(X)$.
\pause
\item Note that if $K_X+\Delta$ is big, then there is an effective divisor
$D\sim _{\mathbb {R}}K_X+\Delta$ and so for all $0<\epsilon \ll 1$
the pair $(X,\Delta+\epsilon D)$ is KLT, $\Delta+\epsilon D$ is big  
and $K_X+\Delta+\epsilon D\sim _{\mathbb {R}}(1+\epsilon)(K_X+\Delta)$.
\pause
\item We would like to find a finite sequence of KLT 
log pairs $(X_i,\Delta _i)$ with $\Delta _i$ big 
and birational maps
$X\dasharrow X_{1}\dasharrow X_{2}$ such that $$\mathfrak {R}(K_{X_i}+\Delta _i
)\cong \mathfrak {R}(K_{X_{i+1}}+\Delta _{i+1}),$$
and $K_{X_t}+\Delta _t$ is nef.
\pause
\item Then by the base point free theorem, $K_{X_t}+\Delta _t$ is
semiample so that $\mathfrak {R}(K_{X_t}+\Delta _t)$ is finitely generated (if $\Delta \in {\rm Div}_\Q$).
\end{itemize}

\end{frame}
\begin{frame}
 \frametitle{The MMP with scaling II}
\begin{itemize}
\item Assume that $A$ is ample and $K_X+\Delta+A$ is nef. 
\pause
\item Let $\lambda ={\rm inf}\{t\geq 0|K_X+\Delta +tA\}$ is nef.
\pause
\item If $\lambda =0$ we STOP ($K_X+\Delta$ is nef). 
\pause
\item Otherwise, by the cone theorem,
there is an extremal ray $R$ such that $(K_X+\Delta +\lambda A)\cdot R=0$ and $(K_X+\Delta +tA)\cdot R<0$
for $0\leq t<\lambda$.
\pause
\item Let ${\rm cont}_R:X\to Z$ be the corresponding morphism which contracts 
a curve $C$ iff $[C]=R$.
\pause
\item If ${\rm cont}_R$ is not birational, we STOP (we
have a Mori-Fano fiber space and $\mathfrak {R}(K_X+\Delta)\cong \mathbb{C}$).
\end{itemize}
\end{frame}
\begin{frame}
\frametitle{The MMP with scaling III}
\begin{itemize}
\item If ${\rm cont}_R$ is birational and divisorial (i.e. it contracts a divisor), then we replace $X$ by the image of the corresponding contraction and we may continue this procedure.
\pause
\item If ${\rm cont}_R$ is birational and small, (i.e. ${\rm Exc}({\rm cont}_R)$ has codimension $\geq 2$) then the image of the corresponding contraction
is not KLT. Therefore, we can NOT replace $X$ by the corresponding contraction!
\pause
\item Instead we replace $(X,\Delta)$ by the corresponding flip 
$\phi :X\dasharrow
X^+$ (assuming it exists).
\pause
\item Let $\Delta ^+=\phi _* \Delta$. The flip is a codimension $2$ surgery 
which replaces $K_X+\Delta$ negative curves by
 $K_{X^+}+\Delta^+$ positive curves.
\end{itemize}
\end{frame}

\begin{frame}
 \frametitle{Flips}
\begin{itemize}
\item More precisely, $f={\rm cont} _R:X\to Z$ is a flipping contraction i.e. a small birational morphism with $\rho (X/Z)=1$ and $-(K_X+\Delta)$ ample over $Z$.
\pause
\item The flip $f^+:X^+\to Z$ (if it exists) is a small birational morphism such that $\rho (X^+/Z)=1$ and $K_{X^+}+\Delta  ^+$ is  ample over $Z$.
\pause
\item The flip $f^+:X^+\to Z$ (if it exists) is given by
$$X^+={\rm Proj}_Z \bigoplus_{m\in \mathbb {N}} f_*\OO _X( [m(K_X+\Delta)]).$$
\pause
\item Note that after peturbing $\Delta$, we may assume that $\Delta \in{\rm Div }_\Q (X)$.
\end{itemize}

\end{frame}
\begin{frame}
 \frametitle{Flips II}
\begin{itemize}
\item The problem of existence of flips is local over $Z$ and hence we may
assume that $Z$ is affine. It suffices then to show that
$\mathfrak{R}(K_X+\Delta)$ is finitely generated.
So this is a local version of the original problem.
\pause
\item To run this MMP with scaling, we must therefore show that such flips
exist and terminate (i.e. there is no infinite sequence of such flips).
\pause
\item Note that after each divisorial contraction, the Picard number drops by $1$ and so there are only finitely many divisorial contractions.
\end{itemize}

\end{frame}

\section{Towards the minimal model program in higher dimensions}
\subsection{Recent results}

\begin{frame}
 \frametitle{Existence and termination of flips}
\begin{thm}[Birkar-Cascini-Hacon-M$^{\rm c}$Kernan] Let $(X,\Delta)$ be a KLT log pair and $f:X\to Z$ a flipping contraction, then the flip $f^+:X^+\to Z$ of $f$ exists.

If moreover $\Delta $ is big, then any sequence of flips for the $K_X+\Delta$ MMP with scaling must terminate.\end{thm}
\pause
\begin{cor} Let $(X,\Delta )$ be a KLT pair with $K_X+\Delta$-pseudo-effective
and $\Delta$ big, then $K_X+\Delta$ has a minimal model. (In particular flips exist.)
\end{cor}
\pause
$K_X+\Delta$-pseudo-effective means that given any ample $\mathbb{Q}$-divisor $H$, $K_X+\Delta+H$ is big (i.e. $\kappa (K_X+\Delta+H)$ $=\dim X$).
This is a necessary condition.

\end{frame}

\begin{frame}
\frametitle{Finite Generation}
\begin{cor} [Birkar-Cascini-Hacon-M$^{\rm c}$Kernan, Siu]
Let $X$ be a smooth projective variety, then
$\mathfrak{R}(K_X)$ is finitely generated.
\end{cor}
\begin{itemize}\pause
\item Proof: Assume that $X$ is of general type.
We may find $\Delta \sim \epsilon K_X$ with $0<\epsilon \ll 1$.
A minimal model for $K_X+\Delta$ is then a minimal model for $K_X$.
The result then follows from the base point free theorem.
\pause
\item As mentioned above, the general case is a consequence of a result of Fujino and Mori.
\pause
\item Siu's proof relies on analytic techniques.
\end{itemize}
\end{frame}
\subsection{The structure of the proof}
\begin{frame}
\frametitle{The structure of the proof}
We proceed by induction on $n=\dim X$. We show that for an $n$-dimensional KLT pair $(X,\Delta)$:

\pause
{\bf Claim 1.} Flips exist.

{\pause
\bf Claim 2.} If $K_X+\Delta \sim _{\mathbb{R}}D\geq 0$ and $\Delta$ is big, then
$(X,\Delta )$ has a minimal model.

\pause
{\bf Claim 3.} The set of minimal models for KLT pairs with $A\leq \Delta
\leq \Delta _0$ where $A$ is ample and $(X,\Delta _0)$ is KLT, is finite.

\pause
{\bf Claim 4.} If $\Delta $ is big and $K_X+\Delta$ is pseudo-effective, then there is a $\Q$-divisor $D$ such that
$K_X+\Delta \sim _{\mathbb {R}}D\geq 0$.

\pause
All statements should be relative over a base.

\end{frame}
\begin{frame}
\frametitle{The structure of the proof II}
\begin{itemize}
\item Roughly speaking, the inductive structure of the proof is:
$$ (1-4)_{n-1}\Rightarrow 1_n\Rightarrow 2_n\Rightarrow 3_n \Rightarrow 4_n.$$
\pause
\item Note that that $(1-4)_{n-1}\Rightarrow 1_n$ was obtained by
Hacon M$^{\rm c}$Kernan in 2005 using ideas of Shokurov and 
an extension result base on work of Siu, Kawamata and Tsuji.

\pause
\item The remaining implications heavily use the ideas of the MMP established by
Kawamata, Koll\'ar, Mori, Reid, Shokurov and others.
\end{itemize}


\end{frame}

\subsection{Existence of Flips}

\begin{frame}
\frametitle{Reduction to PL Flips}
\begin{itemize}
\item Shokurov has shown that to construct flips it suffices to construct PL-flips.
\pause
\item That is, we assume that $\Delta$ has a component $S$ of multiplicity $1$, $-S$ is relatively ample and $K_X+\Delta$ is ``almost KLT''.
\pause
\item By adjunction $(K_X+\Delta )|_S=K_S+\Delta _S$ and $(S,\Delta _S)$ is KLT. 
\pause
\item As we have remarked above, the existence of flips is equivalent to showing that $\mathfrak {R}=\mathfrak {R}(K_X+\Delta )$
is finitely generated. (Here we assume $Z$ is affine.)
\pause
\item This is equivalent to showing that the algebra
$\mathfrak {R}_S$ given by the image of the restriction map
$$\bigoplus _{m\in \mathbb{N}}H^0(\OO _X(m(K_X+\Delta )))\to 
\bigoplus _{m\in \mathbb{N}}H^0(\OO _S(m(K_S+\Delta _S)))$$
is finitely generated.
\end{itemize}
\end{frame}


\begin{frame}
 \frametitle{Existence of flips}
\begin{itemize}
\item If $\mathfrak {R}_S=\oplus _{m\in \mathbb{N}}H^0(\OO _S(m(K_S+\Delta _S)))$ we would be done by induction on the dimension.
\pause
\item This is too much to hope for.
\pause
\item Instead, we show that $\mathfrak {R}_S$ has very special properties.
\pause
\item Roughly speaking we show that there exists a divisor $0\leq \Theta \leq \Delta _S$ on $S$ such that
$$\mathfrak {R}_S=\bigoplus _{m\in \mathbb{N}}H^0(\OO _S(m(K_S+\Theta))).$$
\pause
\item This is proved using some extension theorems that were inspired by work of Siu, Kawamata and Tsuji.
\pause
\item These extension theorems determine/define $\Theta$.
\end{itemize}
\end{frame}
\subsection{Existence of minimal models}

\begin{frame}
\frametitle{Termination of MMP with scaling}
\begin{itemize}
\item Since we are unable to show that an arbitrary sequence of flips 
must terminate, we run a MMP with scaling.
\pause
\item The idea is that if $A$ is sufficiently ample, then $K_X+\Delta +A$ is ample and so it is a minimal model. We then consider $K_X+\Delta +tA$ and let $t$ go to $0$.
\pause
\item Proceeding this way we get a sequence of pairs $(X_i,\Delta _i+\lambda _i A_i)$ such that $K_{X_i}+\Delta _i+\lambda _i A_i$ is nef and trivial on some $K_{X_i}+\Delta _i$ extremal ray. We then flip/contract.
\pause
\item Each $(X_i,\Delta _i+\lambda _i A_i)$ is a minimal model for $(X,\Delta +\lambda _iA)$.
\end{itemize}
\end{frame}
\begin{frame}
\frametitle{Termination of MMP with scaling II}
\begin{itemize}
\item If we could assume Claim $3_n$, we would be done. 
(As $\Delta $ is big, we may assume that $A'\leq \Delta\leq \Delta +tA$ for some ample divisor $A'$ and we let  $\Delta_0=
\Delta +A$ .)
\pause
\item Since we can only assume $3_{n-1}$, we replace $(X,\Delta)$ by
an appropriate DLT pair $(X',\Delta ')$ and we run a MMP with scaling
such that all flips/contractions are contained in $[\Delta ']$.

\pause
\item This step is similar to Shokurov's reduction to PL-flips.
\end{itemize}
\end{frame}
\subsection{Finiteness of models}
\begin{frame}
\frametitle{Finiteness of models}
\begin{thm} Assume that $(X,\Delta _0)$ is KLT, $A$ is an ample $\mathbb {Q}$-divisor and $\Delta =\sum _{i=1}^r
\delta _i \Delta _i$. Let
$\mathcal{C}\subset [0,1]^r$ be a rational polytope such that
for any $(\delta _1,\cdots ,\delta _r)\in \mathcal{C}$, the pair
$(X,\sum _{i=1}^r
\delta _i \Delta _i+A)$ is KLT. 

Then the set of isomorphism classes $$\{Y|Y\ {\rm is\ a\ min.\ model\
for \ }(X,\Delta +A);\ 
(\delta _1,\cdots ,\delta _r)\in \mathcal{C}\}$$
is finite.\end{thm}
\pause
The proof is based on ideas of Shokurov (1996). It requires a compactness argument
and so it is necessary to work with $\mathbb{R}$-divisors.
\end{frame}
\begin{frame}
\frametitle{Finiteness of models II}
\begin{itemize} 
\item We pick $\Delta \in \mathcal{C}$ and work in a sufficiently small neighborhood $U$ of $\Delta$.
\pause
\item After running a MMP with scaling, we may assume that $K_X+\Delta$ is nef.\pause
\item By the base-point free theorem, there is a morphism 
$f:X\to Z$ such that $K_X+\Delta \sim _{\mathbb R}f^*H$.
\pause
\item Working over $Z$, we may assume that $K_X+\Delta \sim _{\mathbb R} 0$.
\pause
\item Therefore $\mathcal{C}\cap U$ is a cone over $\Delta$ and we are 
done by induction on the dimension of the subspace of ${\rm Div }_{\mathbb {R}}(X)$ spanned by the components of $\Delta$.
\end{itemize}

\end{frame}

\subsection{Non-vanishing}
\begin{frame}
\frametitle{Non-vanishing}
\begin{itemize} 
\item We must show that if $\Delta$ is big and $K_X+\Delta$ is 
pseudo-effective, then $K_X+\Delta \sim _{\mathbb {R}}D\geq 0$.
\pause
\item If $h^0(\OO _X(m(K_X+\Delta)+A))$ is bounded (for $A$ ample),
then the result follows from work of Nakayama.
\pause
\item If $h^0(\OO _X(m(K_X+\Delta)+A))$ is un-bounded, then we can produce a
log canonical center not contained in the stable base locus of
$K_X+\Delta+\epsilon A$ for $0<\epsilon \ll 1$.
\pause
\item So we study a log smooth PLT pair $(X',\Delta ')$ where $S=[\Delta ']$
is irreducible.
\pause
\item Then $K_S+(\Delta '-S)|_S$ is pseudo-effective, $(\Delta '-S)|_S$ is 
big and so $h^0(\OO _S(m(K_S+(\Delta '-S)|_S)))>0$ for some $m>0$.
\pause
\item Using Kawamata-Viehweg vanishing, we can lift this section to $m(K_X+\Delta)$.
\end{itemize}

\end{frame}
\subsection{Open problems}
\begin{frame}
\frametitle{Open problems}
\begin{itemize}
\item Abundance: If $K_X+\Delta$ is nef and KLT show that it is semiample.
(When $\Delta$ is big, this follows from the Base Point Free Theorem.
In dim. 2, it follows from the classification. It also holds in dim. 3)
\pause
\item Existence of Minimal models for Log Canonical pairs (should be equivalent to Abundance).
\pause
\item Termination of flips (would follow from the Borisov-Alexeev-Borisov
conjecture).
\pause
\item Characteristic $p$ (need resolution of singularities; there is no Kawamata-Viehweg vanishing Thm, but .... ).
\end{itemize}
\end{frame}

\end{document}
